\chapter{Methodology}
\label{chp:methodology}

\section{Variables}
\begin{itemize}
    \item \textbf{Plaintext modulus} $t$ - Coefficients of plaintext polynomial are taken mod ${t}$
    \item \textbf{Ciphertext modulus} $Q$ - Decomposed into $Q=q_0q_1\dots q_L$
\end{itemize}

\section{Notes}
In Brakerski's scheme, similar to previous FHE schemes based on LWE, the ciphertext $\overset{_\rightarrow}{c}$ and secret key $\overset{_\rightarrow}{s}$ are $n$-dimensional vectors whose dot product $\langle \overset{_\rightarrow}{c}, \overset{_\rightarrow}{s}\rangle \approx \mu$ equals the message $\mu$, up to some small ``error'' that is removed by rounding.

\section{Residue Number System}
In \gls{rns}, each ciphertext is already decomposed mod $Q_L = q_0q_1\dots q_L$. If, then, we select a gadget consisting of $(q_0, q_1, \dots, q_L)$, we get the decomposition ``for free''.

